% chktex-file 45
\section{Introduction}
%%%%%%%%%%%%%%%%%%%%%%%%%%%%%%

A symmetric map $d\: X \times X \to [0,\infty]$ that vanishes on the diagonal and satisfies the triangle inequality is called a \emph{generalized pseudometric}.
A generalized pseudometric that takes no infinite values is called a \emph{pseudometric}.
A generalized pseudometric that vanishes only on the diagonal is called a \emph{generalized metric}, and a pseudometric with this property is called a \emph{metric}.
The Gromov--Hausdorff distance (which we abbreviate to GH) measures how different two metric spaces are.
It was introduced independently by D.~Edwards~\cite{Edwards1975} in 1975 and by M.~Gromov~\cite{Gromov1981, Gromov1981SM} in 1981. This value is the greatest lower bound of the Hausdorff distances between the images of isometric embeddings of the two spaces into all possible ambient spaces.
An equivalent definition via the distortion of correspondences and the properties of this distance are described in detail in~\cite{BurBurIva} (see Section~\ref{sec: preliminaries} below).
The classical Gromov--Hausdorff distance does not take the topology into account: spaces that are topologically different but metrically close can be at a small or even zero distance from each other.

Usually the Gromov--Hausdorff distance is considered on the set of compact metric spaces. We consider it on the class of all metric spaces, so we work in the von~Neumann--Bernays--G\"odel axiom system, where proper classes, which generalize the notion of a set, are allowed in addition to sets.
The proper class of all metric spaces considered up to isometry is denoted by $\GH$; its structure as a proper class is described in detail in~\cite{BogatyiTuzhilinGHClass}.
We treat it as a class of representatives: exactly one space with a fixed metric is chosen from each isometry class of metric spaces.
The notion of a generalized pseudometric is naturally defined on this proper class.

Instead of arbitrary correspondences in the formula of Proposition~\ref{prop: gh via maps}, one can take pairs of maps $f\: X \to Y$ and $g\: Y \to X$ that generate the correspondence $R_{f,g}$.
The \emph{continuous Gromov--Hausdorff distance} $\dGH{C}$ is the modification in which these maps are required to be continuous.
This construction can be generalized: take a subcategory of the category of metric spaces and take the infimum over pairs of its morphisms. The resulting $\FF$-distance $\dGH{\FF}$ ($\FF$ is the class of morphisms of the subcategory) is considered in Section~\ref{sec: f-dist}.
Monotonicity in the class of morphisms (Proposition~\ref{prop:FFmonotone}) implies, in particular, that $\dGH{\FF} \geq \dGH{}$.

The continuous distance was introduced by Lim, M\'emoli and Smith~\cite{LimMemoliSmith} in their comparison of spheres: for the unit spheres $\mathbb{S}^m$ and $\mathbb{S}^n$ with $1 \leq m < n < \infty$ it turned out to be $\pi/2$, while the classical distance is strictly smaller.
Another version, which does not satisfy the triangle inequality, was proposed by Lee and Morales~\cite{LeeMorales2022} for dynamical systems and partial differential equations.
Bogatyi and Tuzhilin~\cite{TuzCont} developed the general theory of the continuous distance (in the Lim--M\'emoli--Smith version): in particular, they proved the triangle inequality and showed that this distance is intrinsic, like the classical one, but incomplete, unlike the classical one.
The generalization to set-valued maps and semicontinuous correspondences is developed in~\cite{SemenovTuzhContGH}.

The difference between the continuous and the classical distances suggests that a small extension of the class of continuous maps can make the corresponding distance equal to the classical GH-distance: for example, one can achieve this by allowing a map to have finitely many points of discontinuity.

As an example, consider the two-point simplex $\Delta_2 = \{a, b\}$ (the only non-zero distance is~1, see the precise definition in Section~\ref{sec: preliminaries}) and the segment $[0,\, 2]$.
Every continuous map $g\: [0,\, 2] \to \Delta_2$ is constant, so $\dis g = \diam[0,\, 2] = 2$, and hence $\dGH{C}(\Delta_2,\, [0,\, 2]) \geq 1$.
If we allow one point of discontinuity, then the corresponding modified distance coincides with the classical one, $\dGH{}(\Delta_2,\, [0,\, 2])$, which is equal to $1/2$ (see~\cite[Thm.~2.12]{GrigorIvanTuzSimplex} for details).
Thus $\dGH{}(\Delta_2,\, [0,\, 2]) = 1/2$, while $\dGH{C}(\Delta_2,\, [0,\, 2]) = 1$.
So for this pair one point of discontinuity is enough to recover the classical distance.
Moreover, for any $|I| \ge 2$ the vertices of the cube ${[0,\, 1]}^I$ and its skeleton have $\dGH{C, \alpha}$ strictly larger than the classical distance for $\alpha = 0$, and, if $I$ is infinite, for all $\alpha < 2^{|I|}$ (see Example~\ref{example: cube skeleton}).

We impose conditions on the maps not on the whole space, but only on a subset of full measure of the support.
First we weaken the continuity requirement in this way: we require the map to be continuous on a subset of full measure of the support that is open in the topology induced on the support.
Then we require the same for local constancy: the map must be locally constant on an open subset of the support of full measure.
This gives the \emph{partially continuous} distance $\dGH{aec}$ and the \emph{partially locally constant} distance $\dGH{aeconst}$ on the class $\GH_{mm}$ of Gromov metric triples $(X,\, d,\, \mu)$, where $\mu$ is an arbitrary boundedly finite Borel measure. The corresponding distances on the class of ordinary metric spaces are defined through them (see Section~\ref{sec: partially smooth}).
The distances $\dGH{aec}$ and $\dGH{aeconst}$ are weaker than $\dGH{C}$, and in this work we show that they coincide with the classical distance $\dGH{}$ (Theorem~\ref{theorem: aec equal original}).

This work also considers restricting the class of continuous maps to smooth ones: for manifolds with a smooth structure it is natural to define the \emph{smooth Gromov--Hausdorff distance} $\dGH{C_k}$ --- the infimum over $C^k$-smooth maps.
It turns out that the smooth structure does not give a finer distinction than the continuous one.

Putting the results together, we obtain the hierarchy
$$
  \dGH{} = \dGH{aec} = \dGH{aeconst} \leq \dGH{C} = \dGH{C_k},
$$
where on the left is the weakest distance, the classical one, which does not take the topology into account; next come the measure-controlled partially continuous and partially locally constant distances, which coincide with it; on the right are the stronger distances that take the topology and smoothness into account, the continuous and the smooth ones, which also coincide with each other.

The author is grateful to his scientific advisor Prof.~A.A.~Tuzhilin and to Prof.~A.O.~Ivanov for formulating the problem and their constant attention to the work.
